\documentclass[11pt]{article}
\usepackage{graphicx}
\usepackage{amssymb}
\usepackage{bm}
\usepackage{mathtools}
\usepackage{amsmath}
\usepackage{commath}
\usepackage{amsthm}
\usepackage{enumitem}
\usepackage{float} 
\usepackage{array}
\usepackage{outlines}
\usepackage{setspace}
\usepackage[colorinlistoftodos]{todonotes}
\setlist[enumerate,1]{label=\alph*)}
\setlist[enumerate,2]{label=\roman*)}
\newtheorem{theorem}{Theorem}
\newtheorem{lemma}{Lemma}
\newtheorem{remark}{Remark}
\newtheorem{definition}{Definition}
\newtheorem{prop}{Proposition}
\newcommand*\diff{\mathop{}\!\mathrm{d}}
\usepackage{tikz}
\usetikzlibrary{arrows,shapes,trees, calc}

%%% *** PREAMBLE *** %%%
\title{A generalised weak convergence result for the mathematical equations of the polluted atmosphere expressed in downwind matching coordinates}

\author{
  \large{\bf{Donatella Donatelli} and \bf{N\'ora Juh\'asz}
  \thanks{The research of DD and NJ leading to these results has received funding from the European Union's Horizon 2020 Research and Innovation Programme under the Marie Sklodowska-Curie Grant Agreement No 642768 (Project Name: ModCompShock).}} \\[1ex]
  \normalsize Department of Information Engineering, Computer Science and Mathematics \\
  \normalsize University of L'Aquila \\
  \normalsize 67100 L'Aquila, Italy. \\
  \normalsize {donatella.donatelli@univaq.it},  {norjuh@univaq.it}
}

\date{}

%%%***  CONTENT *** %%%
\begin{document}

\maketitle

\begin{abstract}

A widely used approach to mathematically describe the atmosphere is to consider it as a geophysical fluid in a shallow domain --- and thus to model it using classical fluid dynamical equations combined with the explicit inclusion of an $\epsilon$ parameter representing the small aspect ratio of the physical domain. In our previous paper we proved a weak convergence theorem for the polluted atmosphere described by the Navier-Stokes equations extended by an advection-diffusion equation. We obtained a justification of the generalised hydrostatic limit model including the pollution effect described in the case of a classical, east-north-upwards oriented local Cartesian coordinates. Here we give an improvement of this statement: we consider a meteorologically more meaningful coordinate system, incorporate the analytical consequences of this coordinate change into the governing equations, and verify that the weak convergence still holds for the altered velocity -- concentration system.

\end{abstract}

{\bf Keywords:} downwind-matching coordinate system, Navier-Stokes equations, shallow domains, pollution evolution equation, hydrostatic approximation, compactness, weak solutions.

\newpage

\tableofcontents
\newpage

%%% *** INTRODUCTION *** %%%
\section{Introduction} \label{Introduction}

A local domain on the surface of the rotating Earth is often described by a Cartesian coordinate system where the $x, y, z$ axes represent the directions towards east, north, and upwards, respectively, as in this specific system the rotation vector, and as a consequence, the Coriolis force takes a simpler form. However, for modelling wind and pollution effects in the atmosphere there is another naturally distinctive viewpoint for fixing the orientation of the coordinate system: the introduction of downwind and crosswind, and to match the horizontal axes according to these meteorological parameters ($\cite{goyal2011mathematical}$). In this paper we take this idea as a foundation, and we migrate the result of our previously stated weak convergence theorem to this more considerately chosen coordinate system.

Here talk a little bit about the original convergence theorem to make the new one more self-contained and for details refer to the first article.

As in this scenario the advection is dominant compared to diffusion effects along the $x$ axis, this approach naturally brings an additional $\epsilon$ term in the Laplacian of the concentration equation, the simplification of the limit model, and brings the challenge of proving a weak convergence result for the product of the vertical velocity and the concentration, as in the energy inequality the $\norm{\partial_1 C^\epsilon}$ term is replaced by $\norm{\sqrt{\epsilon} \partial_1 C^\epsilon}.$ This issue is overcome by the construction of new estimates obtained by multiplying the governing equations by higher order terms.

Modified boundary conditions. Explain the necessity of changing the sediment functions.

%%% *** THE POLLUTED ATMOSPHERE MODEL *** %%%
\section{The polluted atmosphere model} \label{eqsdescribingatmpluspollution}

As we highlighted before, the key point of this new model is to integrate some of the main ideas of a Gaussian plume dispersion model into the system describing the polluted atmosphere. Therefore we begin by fixing the local Cartesian coordinate system to be adjusted to the downwind direction, i.e. $z$ is oriented upwards, while $x$ is the downwind, and $y$ is the crosswind direction. Another assumption the Gaussian dispersion models typically use as well is that advection is more dominant than diffusion along the downwind direction, which leads to the appearance of an additional $\epsilon$ term in the concentration equation's Laplacian. Namely, we have

\begin{equation} \abs{u_1 \partial_{1} C} \gg \abs{ K_x \partial_{11} ^2 C }\end{equation}
    which leads us to using this additional scaling equation
    \[ K_x = \epsilon K_1.\]

\begin{remark} \normalfont
  The inclusion of these two features in a pollution model are legitimate when the following circumstances hold:
  \begin{itemize}
    \item We are considering a time interval which is not too long in the sense that it is reasonable to assume that we have a permanent, relatively stable main wind direction, i.e. fixing the horizontal coordinates is reasonable in the first place,
    \item There is an adequate, relatively high wind speed --- low wind speeds in general lead to a situation where three dimensional diffusion dominates, leading to an error caused by the dominant advection assumption. Fast changing wind direction and low windspeed resulting in fast pollutant settling can make it extremely challenging to obtain reliable results from Gaussian models.
  \end{itemize}
\end{remark}

Choosing this coordinate system results in changes in three main parts of the system:
\begin{itemize}
  \item Describe the Coriolis force and the additional terms in the equations
  \item Describe the boundary terms.
  \item The altered Laplacian in the concentration equation.
\end{itemize}

\begin{remark} \normalfont
The intuitive meaning of what we mean by downwind and the velocity in this new coordinate system. The downwind direction is an approximately steady-in-time direction, it can be seen as the direction marked by the wind direction trackers on an airport or meteorological point. When we consider the downwind, we neglect the small changes in its orientation, it can be viewed as a time-averaged direction as long as the fluctuation is relatively small. On the other hand, the velocity is a three dimensional, time dependent vector that changes second after second. 

\end{remark}
    
    We arrive to the merged anisotropic equations of the polluted atmosphere above inland surface:
  
    \begin{equation} \label{ANSC1}
    \partial_t u^\epsilon _1 + \mathbf{u}^\epsilon \cdot \nabla u^\epsilon _1 -\Delta_\nu u^\epsilon _1 -\alpha u_2^\epsilon + \epsilon \beta u_3^\epsilon + \partial_1 p^\epsilon = 0 \text{ in } \Omega \times (0,T)
    \end{equation}
    
    \begin{equation} \label{ANSC2}
    \partial_t u_2^\epsilon + \mathbf{u}^\epsilon \cdot \nabla u_2^\epsilon -\Delta_\nu u_2^\epsilon + \alpha u_1^\epsilon + \partial_2 p^\epsilon = 0 \text{ in } \Omega \times (0,T)
    \end{equation}
    
    \begin{equation} \label{ANSC3}
    \epsilon^2 \{  \partial_t u_3^\epsilon + \mathbf{u}^\epsilon \cdot \nabla u_3^\epsilon -\Delta_\nu u_3^{\epsilon} \} - \epsilon \beta u_1^\epsilon + \partial_3 p^\epsilon = 0 \text{ in } \Omega \times (0,T)
    \end{equation}
    
    \begin{equation} \label{ANSC4}
    \nabla \cdot \mathbf{u}^\epsilon= 0 \text{ in } \Omega \times (0,T)
    \end{equation}
    
    \begin{equation} \label{ANSC5}
    \partial_t C^\epsilon + \mathbf{u^\epsilon} \cdot \nabla C^\epsilon = \epsilon K_1 \partial_{11} ^ 2 C^\epsilon + K_2 \partial_{22} ^ 2 C^\epsilon + K_3 \partial_{33} ^ 2 C^\epsilon + S_\epsilon \text{ in } \Omega \times (0,T)
    \end{equation}

The initial condition for the velocity and the pollution concentration are
    
    \begin{equation} \label{ANSC7}
    \mathbf{u}^\epsilon (\cdot, t = 0) = \mathbf{u}_0, \quad C^\epsilon (\cdot, t=0) = C_0  \quad \text{ in } \Omega.
    \end{equation}
    


  
  Finally, the boundary conditions to describe the anisotropic system can be summarised as follows:
  
    \begin{equation} \label{ANSC8}
      \begin{split}
        & \mathbf{u}^\epsilon = 0 \text{ on } \Gamma \times (0,T), \\
        &C^\epsilon = 0 \text{ on } ( \Gamma_U \cup \Gamma_L) \times (0,T), \\
        & \partial_2 C^\epsilon = \partial_3 C^\epsilon = 0 \text{ on } \Gamma_G \times (0,T), \quad \partial_1 C^\epsilon \in L^2 L^2 (0,T; \Gamma_G) \\
      \end{split}
    \end{equation}


If we assume $\mathbf{u}^\epsilon = O(1)$, then neglecting the $\epsilon^2$ and $\epsilon$ in the anisotropic equations, we formally arrive to the hydrostatic Navier-Stokes equations (i.e., primitive equations) combined with pollution.

    \begin{equation} \label{HNSC1}    
    \partial_t u_1 + \mathbf{u} \cdot \nabla u_1 - \Delta_\nu u_1 -\alpha u_2 + \partial_1 p = 0 \text{ in } \Omega \times (0,T)
    \end{equation}
    
    \begin{equation} \label{HNSC2}   
    \partial_t u_2 + \mathbf{u} \cdot \nabla u_2 - \Delta_\nu u_2 -\alpha u_1 + \partial_2 p = 0 \text{ in } \Omega \times (0,T)
    \end{equation}
    
    \begin{equation} \label{HNSC3}   
    \partial_3 p = 0 \text{ in } \Omega \times (0,T)
    \end{equation}
    
    \begin{equation} \label{HNSC4}   
    \nabla \cdot \mathbf{u} = 0 \text{ in } \Omega \times (0,T)
    \end{equation}
    
    \begin{equation} \label{HNSC5}   
    \partial_t C + \mathbf{u} \cdot \nabla C = K_2 \partial_{22} ^ 2 C + K_3 \partial_{33} ^2 C + S \text{ in } \Omega \times (0,T)
    \end{equation}
    
    \begin{equation} \label{HNSC6}   
    u_1 = u_2 = u_3 n_3 = 0 \text{ on } \Gamma \times (0,T)
    \end{equation}
    
    \begin{equation} \label{HNSC7}   
    u_i (\cdot, t=0) = u_{0 i}, C (\cdot, t=0) = C_0 \text{ in } \Omega, \text{$i$ = 1,2}
    \end{equation}
    
    \begin{equation} \label{HNSC8}
      \begin{split}
        &C = 0 \text{ on } ( \Gamma_U \cup \Gamma_L) \times (0,T), \\
        & \partial_2 C = \partial_3 C = 0 \text{ on } \Gamma_G \times (0,T), \partial_1 C \in L^2 L^2 (0,T; \Gamma_G) \\
      \end{split}
    \end{equation}

%%% *** WEAK FORMULATION AND MAIN RESULT*** %%%
\section{Weak formulation and Main Result} \label{weakformulationmainresult}
  
%%% *** PROOF *** %%%
\section{Proof of the main theorem} \label{proof}

Estimate for laplacian v, estimate for C.

Apply C in H1.
 

%%% *** CONCLUDING REMARKS *** %%%
\section{Concluding remarks} \label{concludingremarks}

\nocite{*}
\bibliographystyle{abbrv}
\bibliography{the_u_c_eps_model}

\end{document}

%%% add http://elte.prompt.hu/sites/default/files/tananyagok/AtmosphericChemistry/ch10s03.html
